\section{The global cover}\label{sec:cover}

By Proposition~\ref{prop:compact} it remains to prove $\Def\ge0$ on $\Kbox_{\mathrm{red}}$ and to identify the zeros. This is done by
covering $\Kbox$ with finitely many boxes, on each of which one of a small number of sufficient conditions is verified in ball
arithmetic. In this section we state these conditions as lemmas, prove that each implies the inequality, and describe the cover.
The most technical bounds are stated here in outline; their complete statements, with all constants, are in
Appendix~\ref{app:details}.

\subsection{Discharged boxes and the cover theorem}\label{ssec:discharged}

A \emph{box} is a product $B=S\times O$ of eight closed intervals with rational endpoints, where the \emph{shape box}
$S\subset(0,\infty)^2\times\R^3$ and the \emph{offset box} $O\subset\R^3$ are arbitrary. Since $Y(\xi,f+n)=Y(\xi,f)$ for
$n\in\Z^3$, the functions $e$, $D$ and $\Def$ of $\kappa=(\xi,f)$ are defined, and $\Z^3$-periodic in $f$, for all
$f\in\R^3\setminus\Z^3$. We put
\[
  B_{\mathrm{red}}:=\bigl\{(\xi,f)\in B:\ f\notin\Z^3\text{ and }\xi\text{ satisfies \eqref{eq:R1} and \eqref{eq:R2}}\bigr\},
\]
so that $B_{\mathrm{red}}=B\cap\Kbox_{\mathrm{red}}$ for $B\subset\Kbox$; offset boxes that are not contained in $[-\frac12,\frac12]^3$
occur because some rules are verified in translated coordinates (Section~\ref{ssec:stages}). The boxes that the cover has to
exhaust are contained in $\Kbox$, whose shape part is
$[\frac58,\frac{37}{32}]\times[\frac12,\frac{11}8]\times[0,\frac12]^2\times[-\frac12,\frac12]$. Boxes that carry a rule may extend
beyond it; this is harmless, since Definition~\ref{def:discharged}, Lemma~\ref{lem:discharged} and the rules below hold for all
boxes in the sense just defined, and Theorem~\ref{thm:cover} concerns only points of $\Kbox$. Let
$\iota(r_1,r_2,x,y,z,f):=(r_1,r_2,x,y,z,-f)$; recall that $\Def\circ\iota=\Def$.
The numbers $\eB,\eF$ are known through enclosures $[\eB^-,\eB^+]$, $[\eF^-,\eF^+]$ of width below $10^{-45}$
(Section~\ref{ssec:constcert}); we put $\gap^+:=\eF^+-\eB^-\ge\gap$.

\begin{definition}\label{def:discharged}
A box $B$ is \emph{discharged} if one of the following holds.
\begin{itemize}
\item[(N)] $B_{\mathrm{red}}=\emptyset$.
\item[(E)] $e(\kappa)>\eF$ for all $\kappa\in B_{\mathrm{red}}$.
\item[(D)] $\Def(\kappa)>0$ for all $\kappa\in B_{\mathrm{red}}$.
\item[(A)] There are $f_0\in\frac12\Z^3\setminus\Z^3$, $n\in\Z^3$ and $c>0$ such that, for all $\kappa=(\xi,f)\in B_{\mathrm{red}}$,
\begin{equation}\label{eq:anchorconcl}
  \Def(\xi,f)\ \ge\ \Def(\xi,f_0)+c\,\bigl|R(\xi)(f-f_0-n)\bigr|^2 .
\end{equation}
\end{itemize}
\end{definition}

\begin{lemma}\label{lem:discharged}
\textup{(i)} If $B$ is discharged, then $\Def\ge0$ on $B_{\mathrm{red}}$, and $\Def(\kappa)=0$ for some $\kappa\in B_{\mathrm{red}}$ only if
$X(\kappa)$ is a BCC or an FCC lattice.
\textup{(ii)} If $B$ is discharged, so are $\iota(B)$, every translate $B+(0,n)$, $n\in\Z^3$, and every box contained in $B$.
\textup{(iii)} Let $B=B^\circ\cup\Sigma_1\cup\dots\cup\Sigma_r$ with boxes $B^\circ$ and $\Sigma_i$ such that every $\Sigma_i$ satisfies
\textup{(N)}. If $B^\circ$ satisfies one of \textup{(N)}, \textup{(E)}, \textup{(D)}, \textup{(A)}, then so does $B$.
\end{lemma}

\begin{proof}
(i) (N) and (D): nothing to prove. (E): $\Def=e-\eB-\gap D\ge e-\eB-\gap=e-\eF>0$, since $D\le1$.
(A): $X(\xi,f_0)$ is a lattice, because $2f_0\in\Z^3$, so $\Def(\xi,f_0)\ge0$ by Theorem~\ref{thm:bravais}, and \eqref{eq:anchorconcl} gives
$\Def(\xi,f)\ge0$. If $\Def(\xi,f)=0$ then $R(\xi)(f-f_0-n)=0$, so $X(\xi,f)=X(\xi,f_0)$ is a lattice with $\Def=0$, hence a BCC or
an FCC lattice by Theorem~\ref{thm:bravais}.
(ii) Each condition of Definition~\ref{def:discharged} is a statement about the individual points of $B_{\mathrm{red}}$ (in (A) with
fixed $f_0$, $n$ and $c$), so it is inherited by every box contained in $B$. The maps $\iota$ and $(\xi,f)\mapsto(\xi,f+n)$ preserve
$e$ and $D$, and they map $B_{\mathrm{red}}$ onto $\iota(B)_{\mathrm{red}}$ and onto $(B+(0,n))_{\mathrm{red}}$. This gives the claim
for (N), (E) and (D). In (A), replace $(f_0,n)$ by $(f_0,-n-2f_0)$ for $\iota$, and $n$ by $n$ plus the translation vector for a
translate.
(iii) $B_{\mathrm{red}}\subset B^\circ_{\mathrm{red}}$.
\end{proof}

Note that in (A) the point $(\xi,f_0)$ need not belong to $B$ or to $\Kbox_{\mathrm{red}}$: Theorem~\ref{thm:bravais} holds for every
lattice, however it is represented. Conditions (E) and (D) can never hold on a box that contains a zero of $\Def$; near the zeros,
and near the lattices on which $\Def$ is small, the cover relies on (A). In the rules below, (E) is verified in the form
$e>\eF^+$ and (D) in the form $e-\eB^+-\gap^+\bar D>0$ with an upper bound $\bar D\ge D$; these imply (E) and (D) because
$\gap\le\gap^+$ and $D\ge0$. Part (iii) is used as follows: slabs $\{\kappa\in B:\kappa_i\le a\}$ or $\{\kappa_i\ge a\}$ cut from a box
along a shape coordinate are removed when they satisfy (N), and a rule is applied to the remaining box.

\begin{theorem}[Cover]\label{thm:cover}
There is a finite family $\mathcal F$ of discharged boxes such that for every $\kappa\in\Kbox_{\mathrm{red}}$ at least one of the
points $\kappa$, $\iota(\kappa)$ lies in a box of $\mathcal F$.
\end{theorem}

The family $\mathcal F$, with the data that prove each box discharged, is the certificate described in
Sections~\ref{ssec:stages} and~\ref{sec:computation}. Sections~\ref{ssec:tools}--\ref{ssec:anchor} contain the sufficient conditions.

\subsection{Three tools}\label{ssec:tools}

\emph{The screened series.} Applying Proposition~\ref{prop:ewald} with $\alpha\sqrt\pi/\ltwo$ in place of $\sqrt\pi/\ltwo$ gives,
for every $\alpha>0$,
\begin{equation}\label{eq:screened}
  2\ltwo\,e(\kappa)=\sum_{m\ne0}F_{\mathrm r}(q_m)+\sum_{m}F_{\mathrm r}(q^\times_m)+\sum_{m\ne0}\bigl(1+\cos2\pi m\cdot f\bigr)F_{\mathrm d}(q^*_m)
     -2\alpha-\frac2{\alpha^2},
\end{equation}
with $F_{\mathrm r}(q)=\alpha\,\Ga(\frac12,\pi\alpha^2q)$, $F_{\mathrm d}(q)=\alpha^{-2}\Ga(1,\pi q/\alpha^2)$ and the squared lengths
\begin{align*}
  q_m&=|Rm|^2=r_1^2(m_1+xm_2+ym_3)^2+r_2^2(m_2+zm_3)^2+r_3^2m_3^2,\qquad q^\times_m=|R(m+f)|^2,\\
  q^*_m&=|R^{-\T}m|^2=r_1^{-2}m_1^2+r_2^{-2}(m_2-xm_1)^2+r_1^2r_2^2\bigl(m_3-zm_2+(xz-y)m_1\bigr)^2 .
\end{align*}
For $\alpha=1$ this is \eqref{eq:H2}. All terms are nonnegative, and $F_{\mathrm r}$, $F_{\mathrm d}$ are positive, decreasing and
logarithmically convex; moreover each of them is the Laplace transform of a positive measure. Different rules use different values of
$\alpha$ between $1$ and $\frac65$; a larger $\alpha$ moves weight from the direct to the reciprocal sums.

\emph{Convex supports.} Let $F\in\{F_{\mathrm r},F_{\mathrm d}\}$ and $Q>0$. Then for all $q>0$
\begin{equation}\label{eq:tangent}
  F(q)\ \ge\ F(Q)\exp\bigl[-s\,(q-Q)\bigr]\ \ge\ F(Q)+F'(Q)\,(q-Q),\qquad s:=-F'(Q)/F(Q)>0 .
\end{equation}
For the phase factor $w(t):=1+\cos2\pi t=2\cos^2\pi t$ we use three lower bounds. If an interval $[t_L,t_R]$ contains no point of
$\Z+\frac12$, then $\log w$ is concave on it and lies above its chord:
\begin{equation}\label{eq:secant}
\begin{gathered}
  w(t)\ \ge\ p\,e^{b(t-t_0)}\qquad(t_L\le t\le t_R),\\
  p=\sqrt{w(t_L)w(t_R)},\qquad b=\frac{\log w(t_R)-\log w(t_L)}{t_R-t_L},\qquad t_0=\frac{t_L+t_R}2 ,
\end{gathered}
\end{equation}
where $b:=0$ if $t_L=t_R$, in which case \eqref{eq:secant} is an equality. If $|t-(k+\frac12)|\le d\le\frac12$ for an integer $k$,
then, by the concavity of the sine,
\begin{equation}\label{eq:extinction}
  w(t)=2\sin^2\pi\bigl(t-k-\tfrac12\bigr)\ \ge\ c_d\,\bigl(t-k-\tfrac12\bigr)^2,\qquad c_d=2\Bigl(\frac{\sin\pi d}d\Bigr)^2\quad(d>0),
  \qquad c_0=2\pi^2 ;
\end{equation}
for $d=0$ both sides vanish.
Finally, $|1+e^{iu}+e^{-iv}|^2\ge0$ gives $\cos u+\cos v+\cos(u+v)\ge-\frac32$, so that
\begin{equation}\label{eq:triad}
  w(m_1\cdot f)+w(m_2\cdot f)+w(m_3\cdot f)\ \ge\ \tfrac32\qquad\text{whenever }m_1\pm m_2\pm m_3=0 .
\end{equation}

\emph{The vertex principle.} A function on a box is \emph{separately convex} if it is convex in each coordinate when the other seven
are fixed. Such a function attains its maximum over the box at one of the $2^8=256$ vertices (maximize over one coordinate at a
time), and a separately concave function attains its minimum at a vertex.

\begin{lemma}\label{lem:sepconv}
For fixed $m$, the functions $q_m$, $q^\times_m$ and $q^*_m$ of $\kappa=(r_1,r_2,x,y,z,f)$ are separately convex on
$(0,\infty)^2\times\R^6$, in particular on every box. So is
every linear combination of them with nonnegative coefficients plus an affine function of $f$, and so is $\Psi\circ q$ for every
increasing convex $\Psi$ and $q\in\{q_m,q_m^\times,q_m^*\}$.
\end{lemma}

\begin{proof}
As functions of $r_1$ or of $r_2$ alone they are combinations of $r^2$ and $r^{-2}$ with nonnegative coefficients; as functions of
$x$, $y$, $z$ or $f_i$ alone they are sums of squares of affine functions with nonnegative coefficients.
\end{proof}

No joint convexity is claimed or needed. It is essential, however, that in every application the extremum is taken \emph{after} all
terms have been added: the sum is evaluated at each of the $256$ vertices, and then the extremum over the vertices is taken. This
retains the correlations between the terms, which separate bounds for the individual terms lose.

\subsection{Lower bounds for the energy on a box}\label{ssec:energyrules}

Fix a box $B$, a value of $\alpha$ and a finite set $J$ of terms of \eqref{eq:screened}; the omitted terms are nonnegative. A term
$j\in J$ has a kernel $F_j$, a squared length $q_j(\kappa)$ and, for reciprocal terms, a phase $t_j(\kappa)=m_j\cdot f$, which is affine
in $f$.

\begin{lemma}[Energy rules]\label{lem:energy}
Each of the following implies \textup{(E)} on $B$ (on all of $B$, reduced or not, except in case \textup{(iv)}).
\begin{enumerate}
\item[\textup{(i)}] \emph{Partial sum.} With upper bounds $\bar q_j\ge\sup_Bq_j$ and numbers $0\le\omega_j\le\inf_Bw(t_j)$ (and
$\omega_j=1$ for direct terms),
$\sum_{j\in J}\omega_jF_j(\bar q_j)-2\alpha-2\alpha^{-2}>2\ltwo\eF^+$.
\item[\textup{(ii)}] \emph{Tangent support.} For each $j\in J$ let $\tau_j$ be one of the following functions on $B$:
\begin{itemize}
\item[(a)] $\omega_j\bigl[F_j(Q_j)+F_j'(Q_j)(q_j-Q_j)\bigr]$ with an arbitrary $Q_j>0$ and $\omega_j$ as in \textup{(i)};
\item[(b)] for a reciprocal term whose phase interval over $B$ contains no point of $\Z+\frac12$, the tangent plane at an arbitrary
point $(Q_j,T_j)\in(0,\infty)\times\R$ of the convex function $(q,t)\mapsto F_j(q)\,p_je^{b_j(t-t_{0,j})}$, evaluated at $(q_j,t_j)$,
with $p_j,b_j,t_{0,j}$ from \eqref{eq:secant};
\item[(c)] for a reciprocal term whose phase stays within $d\le\frac12$ of some $k+\frac12$, $k\in\Z$,
$2u_j\zeta_j-u_j^2\,/\bigl(c_dF_j(q_j)\bigr)$ with $\zeta_j=t_j-k-\frac12$ and an arbitrary $u_j\in\R$.
\end{itemize}
Then $\mathcal T:=\sum_{j\in J}\tau_j$ is separately concave, and the condition is
$\min_{v}\mathcal T(v)-2\alpha-2\alpha^{-2}>2\ltwo\eF^+$, the minimum being over the vertices $v$ of $B$.
\item[\textup{(iii)}] \emph{Exponential support.} Let $Q_j>0$ and $s_j=-F_j'(Q_j)/F_j(Q_j)$. Let $W_j=F_j(Q_j)$, $b_j=0$ for direct terms;
for a reciprocal term let either $W_j=p_jF_j(Q_j)$ with $p_j$, $b_j$, $t_{0,j}$ as in \textup{(b)}, or $W_j=\omega_jF_j(Q_j)$, $b_j=0$
with $\omega_j$ as in \textup{(i)}. With $M=\sum_jW_j$, which is required to be positive, and
\[
  U(\kappa):=\frac1M\sum_{j\in J}W_j\bigl[-s_j(q_j(\kappa)-Q_j)+b_j(t_j(\kappa)-t_{0,j})\bigr],
\]
the condition is $M\exp\bigl(\min_vU(v)\bigr)-2\alpha-2\alpha^{-2}>2\ltwo\eF^+$.
\item[\textup{(iii$'$)}] \emph{Phase reserve.} In \textup{(ii)} and in \textup{(iii)}, let $T$ range over triples of three distinct
reciprocal terms of $J$ with $m_1\pm m_2\pm m_3=0$, each of which is bounded with a constant phase factor $\omega_j$:
in \textup{(ii)} by \textup{(a)}, in \textup{(iii)} by the alternative $W_j=\omega_jF_j(Q_j)$, $b_j=0$. Let $s_T\ge0$ satisfy
$\sum_{T\ni j}s_T\le\inf_BF_j(q_j)$ for every $j$. Then $\varrho:=\sum_Ts_T\bigl(\tfrac32-\sum_{j\in T}\omega_j\bigr)_+$ may be added to the
left side of the condition. Terms bounded by \textup{(b)} or \textup{(c)}, or by the chord in \textup{(iii)}, do not qualify.
\item[\textup{(iv)}] \emph{Reduced variant.} In \textup{(ii)(a)} or \textup{(iii)}, let $\Lambda=\sum_j\lambda_jq_j$ be the combination
of squared lengths that occurs with a negative sign, let $\lambda_{e_1}$ be the coefficient of $q_{e_1}=r_1^2$, and let $c$ be one of
$c_1=r_2^2-\frac34r_1^2$, $c_2=r_2^2+r_1^2x^2-r_1^2$. For $0\le\lambda\le\lambda_{e_1}$ replace $\Lambda$ by $\Lambda+\lambda c$
before taking the extremum over the vertices. Then the conclusion \textup{(E)} holds on $B_{\mathrm{red}}$.
\end{enumerate}
\end{lemma}

\begin{proof}
(i) $F_j$ is decreasing and all omitted terms are nonnegative.
(ii) Each $\tau_j$ is a lower bound for the $j$-th term on $B$: (a) by \eqref{eq:tangent}, with $w(t_j)\ge\omega_j\ge0$; (b) because
$w(t)F_j(q)\ge F_j(q)p_je^{b_j(t-t_{0,j})}$ by \eqref{eq:secant}, and the right side is logarithmically convex, hence convex, in $(q,t)$,
so it lies above its tangent planes, also at points $(Q_j,T_j)$ that are not attained on $B$; (c) by \eqref{eq:extinction},
$w(t_j)F_j(q_j)\ge\zeta_j^2/\mathsf R$ with $\mathsf R=1/(c_dF_j(q_j))$, and $\zeta^2/\mathsf R\ge2u\zeta-u^2\mathsf R$. In (a) and (b)
the coefficient of $q_j$ is negative, and in (c) $1/F_j$ is increasing and convex, since $1/F_{\mathrm d}(q)$ is a multiple of
$qe^{\pi q/\alpha^2}$; so each $\tau_j$ is separately concave by Lemma~\ref{lem:sepconv}.
(iii) By \eqref{eq:tangent} and \eqref{eq:secant} the $j$-th term is at least $W_j\exp(u_j)$ with $u_j$ the bracket in the definition
of $U$, and Jensen's inequality gives $\sum_jW_je^{u_j}\ge Me^U$; $U$ is separately concave.
(iii$'$) For $0\le a_j\le F_j(q_j)$ we have $w(t_j)F_j(q_j)\ge\omega_jF_j(q_j)+a_j(w(t_j)-\omega_j)$, the difference being
$(F_j(q_j)-a_j)(w(t_j)-\omega_j)\ge0$. For a term of a triple, $\omega_jF_j(q_j)$ is at least the bound for the $j$-th term that
\textup{(a)}, respectively the alternative $W_j=\omega_jF_j(Q_j)$ of \textup{(iii)}, uses, by \eqref{eq:tangent}. With
$a_j=\sum_{T\ni j}s_T$,
$\sum_ja_j(w(t_j)-\omega_j)=\sum_Ts_T\sum_{j\in T}(w(t_j)-\omega_j)\ge\sum_Ts_T(\frac32-\sum_{j\in T}\omega_j)_+$ by \eqref{eq:triad} and
$w(t_j)\ge\omega_j$. Hence the sum of the terms of $J$ is at least $\mathcal T+\varrho$ in \textup{(ii)}, and at least
$\sum_jW_je^{u_j}+\varrho\ge Me^U+\varrho$ in \textup{(iii)}.
(iv) On reduced shapes $c_2=|Re_2|^2-|Re_1|^2\ge0$ and $c_1\ge0$ by \eqref{eq:R1}, the latter with $m=e_2-ke_1$ for the integer $k$
nearest to $x$, since $|R(e_2-ke_1)|^2=r_1^2(x-k)^2+r_2^2$ and $(x-k)^2\le\frac14$. Hence $\Lambda\le\Lambda+\lambda c$ on
$B_{\mathrm{red}}$, and $\Lambda+\lambda c$ is still separately convex: the coefficient of $r_1^2$ in it is at least
$\lambda_{e_1}-\lambda\ge0$.
\end{proof}

The base points $Q_j,T_j,u_j$ and the weights $s_T$ of the triples are free parameters, and their choice affects the strength of
the bound but not its validity. The base points are proposed by floating-point optimization and determined by exact data stored
in the certificate; the weights are allocated by the verifying program (Section~\ref{ssec:stages}). In the last stages the base
points $(Q_j,T_j)$ are averages of the values of $(q_j,t_j)$ over the
vertices of $B$ with respect to an optimized probability distribution on the vertices. A support that is valid on a box is valid
on every sub-box, and re-evaluating $\min_v\mathcal T(v)$ over the vertices of a sub-box requires no new evaluation of special
functions; this makes deep subdivision cheap. Rules (ii)--(iv) may be applied also to boxes $B$ that meet $\{f\in\Z^3\}$: such
points do not belong to $B_{\mathrm{red}}$, and the functions $\mathcal T$ and $U$ are finite and separately concave on all of $B$.

\subsection{Planar expansion}\label{ssec:planar}

When the cell is elongated or the offset is close to a mid-plane, the three-dimensional sums converge slowly in the bounds above, and
it is better to sum over planes. Let
\[
  A_\parallel=\begin{pmatrix}r_1&r_1x\\0&r_2\end{pmatrix},\qquad a=r_1r_2,\qquad c=\frac1a,\qquad \sigma=(y-xz,\,z),\qquad f=(f_\parallel,f_3),
\]
so that $Y(\kappa)$ consists of translates of the planar lattice $A_\parallel\Z^2$ at the heights $kc$ and $(k+f_3)c$, $k\in\Z$, shifted
in the plane by $k\sigma$ and by $f_\parallel+(k+f_3)\sigma$ (in units of the planar basis). For $p\in\Z^2\setminus\{0\}$ let
$\varkappa_p:=|A_\parallel^{-\T}p|$ and $T_p:=2\pi c\varkappa_p$.

\begin{lemma}[Planar expansion]\label{lem:planar}
For $0<f_3<1$,
\begin{multline}\label{eq:planar}
  2\ltwo\,e(\kappa)=\mathcal M(A_\parallel)+2\pi c^2\bigl(f_3^2-f_3+\tfrac13\bigr)\\
   +\frac1a\sum_{p\ne0}\frac1{\varkappa_p}\Bigl[\sum_{k\ne0}e^{-T_p|k|}\cos(2\pi k\,p\cdot\sigma)
   +\sum_{k\in\Z}e^{-T_p|k+f_3|}\cos\bigl(2\pi p\cdot(f_\parallel+(k+f_3)\sigma)\bigr)\Bigr],
\end{multline}
where, for a real $2\times2$ matrix $A$ with $\det A=a>0$,
\begin{equation}\label{eq:M2}
  \mathcal M(A)=\int_0^\infty t^{-1/2}\Bigl[\sum_{n\in\Z^2}e^{-\pi t|An|^2}-1-\frac1{at}\Bigr]dt
  =a^{-1/2}\Bigl[-4+\sum_{n\ne0}\Bigl(\Ga\bigl(\tfrac12,\tfrac\pi a|An|^2\bigr)+\Ga\bigl(\tfrac12,\pi a|A^{-\T}n|^2\bigr)\Bigr)\Bigr].
\end{equation}
\end{lemma}

\begin{proof}
Write the left side as the Mellin integral $\int_0^\infty t^{-1/2}[\,\cdot\,]\,dt$ of the theta series of $Y$, regularized as in the
proof of Proposition~\ref{prop:ewald}, and apply the two-dimensional Poisson summation formula in each plane. The mode $p=0$ gives the
one-dimensional energy of uniformly charged planes in a neutralizing background, which is the quadratic polynomial in $f_3$; the modes
$p\ne0$ give the exponentially decaying terms through
$\int_0^\infty t^{-3/2}e^{-\pi\varkappa^2/t-\pi c^2s^2t}\,dt=\varkappa^{-1}e^{-2\pi c\varkappa|s|}$. The in-plane term of the first
coset is $\mathcal M(A_\parallel)$; the second expression for it follows by splitting the integral at $t=1/a$ and applying Poisson
summation to the part $t<1/a$. The details are given in Appendix~\ref{app:details}.
\end{proof}

The series \eqref{eq:planar} converges exponentially, with rate $2\pi c\varkappa_p\min(f_3,1-f_3)$ for the cross-coset layers, and it
has few significant terms when $c\varkappa_p$ is not small. Its terms have no sign, so truncation costs explicit tails, which are
geometric series. Two rules are built on it.

\begin{lemma}[Planar energy rules]\label{lem:planarrules}
Let $B=S\times O$ be a box, translated by an integer multiple of $e_3$ if necessary; this is permitted by
Lemma~\ref{lem:discharged}\textup{(ii)}.

\textup{(i)} Let $\eta\le f_3\le1-\eta$ on $B$, with $\eta>0$. By Montgomery's theorem \cite{Montgomery1988} the planar theta function
at fixed area is minimized by the triangular lattice, so $\mathcal M(A)\ge a^{-1/2}M_\triangle$, where $M_\triangle=-3.9210\ldots$ is
the value of the bracket in \eqref{eq:M2} for the triangular lattice; moreover the bracket is bounded below by each of its partial
sums minus $4$, with the squared lengths replaced by upper bounds over $B$. Let $\underline{\mathcal M}$ be a lower bound for
$\mathcal M(A_\parallel)$ of one of these two kinds. If a lower bound over $B$ for the sum of $\underline{\mathcal M}$, of the
zero-mode term and of the terms of \eqref{eq:planar} with $\max|p_i|\le2$ from the two nearest layers of each coset, minus an upper
bound for the absolute values of the omitted terms \textup{(}Section~\ref{ssec:det-fibre}\textup{)}, exceeds $2\ltwo\eF^+$, then
\textup{(E)} holds on $B$.

\textup{(ii)} Let $|f_3-\frac12|\le\eta<\frac12$ on $B$. There is an explicit lower bound for the right side of \eqref{eq:planar} that
depends on the shape box and on $\eta$ only and holds for \emph{all} $f_\parallel\in\R^2$
(Lemma~\ref{lem:det-fibre}); if it exceeds $2\ltwo\eF^+$, then \textup{(E)} holds on $S\times\R^2\times\{|f_3-\frac12|\le\eta\}$.
\end{lemma}

In (ii) the two adjacent cross-coset layers are combined first. For the three shortest waves $p_1=(1,0)$, $p_2=(0,1)$,
$p_3=(1,1)$, whose phases $\theta_i=2\pi p_i\cdot(f_\parallel+(f_3-\frac12)\sigma)$ satisfy $\theta_3=\theta_1+\theta_2$, and positive
weights $w_i$, one has the identity
\[
  \sum_{i=1}^3w_i\cos\theta_i=\tfrac12\bigl|A_0+A_1e^{i\theta_1}+A_2e^{-i\theta_2}\bigr|^2-\tfrac12\bigl(A_0^2+A_1^2+A_2^2\bigr),
  \qquad A_0A_1=w_1,\ A_0A_2=w_2,\ A_1A_2=w_3 .
\]
The nonnegative square controls the contributions of the next waves and of the terms that are odd in $f_3-\frac12$, and the zero-mode
term $2\pi c^2(f_3-\frac12)^2$ controls the dependence on the height; completing a square in the two variables
$|A_0+A_1e^{i\theta_1}+A_2e^{-i\theta_2}|$ and $|f_3-\frac12|$ gives a bound in which $f_\parallel$ no longer appears. The weights
are the coefficients of the two combined layers; the rule requires them to be positive on the shape box for the waves $p_i$ and
$2p_i$, the identity being used for both triples. In this rule $\mathcal M(A_\parallel)$ is bounded below by a finite partial sum of
the positive series in \eqref{eq:M2}, without Montgomery's theorem.

\subsection{The defect on a box}\label{ssec:Dbox}

In this section a \emph{registration} is a triple $(B',f_0,n)$ of a basis matrix $B'$, with $\det B'=1$, of a sublattice of index two
of $\Bz$, its half class $f_0\in\frac12\Z^3\setminus\Z^3$ (so that $B'f_0\in\Bz\setminus B'\Z^3$) and $n\in\Z^3$. For $\kappa=(\xi,f)$ put
\begin{equation}\label{eq:regdata}
\begin{gathered}
  A:=R(\xi)B'^{-1},\qquad C:=A^\T A,\qquad \delta:=f-f_0-n,\qquad h:=R(\xi)\delta,\\
  d_0:=100\,\|C^{1/2}-I\|_\Fro^2,\qquad q:=d_0+200\,|h|^2 .
\end{gathered}
\end{equation}
Then $Y(\kappa)=A\bigl(B'\Z^3\cup(B'\Z^3+B'f_0+A^{-1}h)\bigr)$, and Theorem~\ref{thm:registration} says that $D(\kappa)<1$ if and only if
some registration has $q+\Phi<1$, where $\Phi$ is the sum over unmatched neighbours in \eqref{eq:Dformula}, and then $D=q+\Phi$ for
this registration.

\begin{lemma}[Bounds for the defect]\label{lem:Dbox}
Let $(B',f_0,n)$ be a registration, with the notation \eqref{eq:regdata}.
\begin{enumerate}
\item[\textup{(i)}] If the eigenvalues of $C$ lie in $[\mu_-,\mu_+]$ with $\mu_-\ge\frac{169}{200}$,
$\sqrt{\mu_+}\,(1+|B'\delta|)<\frac{13}{10}$ and $\sqrt{2\mu_-}-\sqrt{\mu_+}\,|B'\delta|>\frac{13}{10}$, then
$D\le28\,\|C-I\|_\Fro^2+200\,|h|^2$.
\item[\textup{(ii)}] If $q<1$, then every neighbour vector that is not matched by the registration is the image of one of the twelve
vectors of $\Bz$ of length $\sqrt2$ and is longer than $\frac65$, and none occurs if
$q\le d_\sharp:=\frac{600}{11}(\sqrt2-\frac{13}{10})^2=0.7115\ldots$ Hence $D=\min\{1,q+\Phi\}$ with $\Phi$ a sum of at most twelve terms
$\phi(|y|)$.
\item[\textup{(iii)}] For every $U\in\SO(3)$ we have $q\le100\,\|A-U\|_\Fro^2+200\,|h|^2$. If $q<1$, then
$D\le100\,\|A-U\|_\Fro^2+200\,|h|^2+\Phi$ with $\Phi$ as in \textup{(ii)}; if $q\ge1$, then $D\le1\le100\,\|A-U\|_\Fro^2+200\,|h|^2$.
\item[\textup{(iv)}] \emph{(Saturation.)} For a shape box $S$ there is a finite list of pairs $(B',f_0)$, computable in exact
arithmetic, that contains, for every $\xi\in S$, all registrations with $d_0<1$ up to rotations preserving $\Bz$. If for every pair of
the list and every $n$ with $|h|<\frac1{14}$ somewhere on $B$ a lower bound for $q+\Phi$ over $B$ is at least $1$, then $D\equiv1$ on $B$.
\item[\textup{(v)}] \emph{(Lattices.)} If $\delta=0$, then $D(\xi,f_0)\le28\,\|C-I\|_\Fro^2$, and $D(\xi,f_0)=d_0$ whenever $d_0<1$.
\end{enumerate}
\end{lemma}

\begin{proof}
(i) is Lemma~\ref{lem:Dupper}(i) with $(\sqrt c-1)^2\le(c-1)^2/(1+\sqrt{\mu_-})^2$ and $100/(1+\sqrt{169/200})^2<28$. (ii) As in Step~6
of the proof of Theorem~\ref{thm:registration}, a vector $Ab+\chi(b)h$ with $|b|=t$ has length at least
$t-\sqrt q\,(t^2/150+1/200)^{1/2}$, where we used that the least singular value $s$ of $A$ satisfies
$\|C^{1/2}-I\|_\Fro^2\ge(1-s)^2+2(s^{-1/2}-1)^2\ge\frac32(1-s)^2$ because $\det A=1$; this exceeds $\frac65$ for $t=\sqrt2$ and
$\frac{13}{10}$ for the next length $t=\sqrt{11}/2$ when $q<1$, and it exceeds $\frac{13}{10}$ for $t=\sqrt2$ when $q\le d_\sharp$. The
trial given by the registration and the polar rotation of $A$ has cost $q+\Phi$, so $D\le\min\{1,q+\Phi\}$, with equality if
$q+\Phi<1$ by Theorem~\ref{thm:registration}; and if $q+\Phi\ge1$, a trial of cost less than $1$ would match the same fourteen
vectors, by Theorem~\ref{thm:registration}(a) (the vectors $Ae+\chi(e)h$, $e\in\Star$, are shorter than
$\frac{11}{10}+\frac1{14}<\frac65$ when $q<1$), and would have the same least cost. (iii) The first claim holds because
$d_0=100\dist_\Fro(A,\SO(3))^2\le100\|A-U\|_\Fro^2$. If $q<1$, the bound is the cost of the trial $\sigma(e)=Ae+\chi(e)h$, $O=U$.
By the proof of (ii), these fourteen vectors are elements of $\mathcal N$ shorter than $\frac65$, and the neighbours that the trial
leaves unmatched are those counted in $\Phi$; the vectors are distinct because $|A(e-e')|\ge\frac9{10}\cdot\frac{\sqrt3}2>|h|$ for
$e\ne e'$. Since eight elements of $\Star$ have $\chi=1$ and $\sum_e\chi(e)e=0$, the matched part of the cost is
$25\,(4\|A-U\|_\Fro^2+8|h|^2)$, as in Steps~4 and~5 of the proof of Theorem~\ref{thm:registration}. If $q\ge1$, the bound is
$D\le1$. (iv) By Theorem~\ref{thm:registration},
$D<1$ at $\kappa$ implies a registration with $q+\Phi=D<1$; then $\|C^{1/2}-I\|_\Fro<\frac1{10}$ gives
$\frac9{10}|B'm|\le|R(\xi)m|\le\frac{11}{10}|B'm|$ for $m\in\Z^3$, which confines the columns of $B'$ to finitely many vectors of $\Bz$
once $\xi$ ranges over $S$, and $|h|^2<\frac1{200}$ confines $n$. (v) Let $\delta=0$. If $d_0<1$, then
$(1-s)^2+2(s^{-1/2}-1)^2<\frac1{100}$ gives $s>\frac{23}{25}$, so the images of the vectors of length $\sqrt2$ have length at least
$s\sqrt2>\frac{13}{10}$, $\Phi=0$ and $D(\xi,f_0)=d_0$ by (ii). For the first claim we may assume $28\|C-I\|_\Fro^2<1$; then
$\|C-I\|_\op<\frac{19}{100}$, the eigenvalues of $C$ exceed $(\frac9{10})^2$, and
$d_0\le\frac{100}{(19/10)^2}\|C-I\|_\Fro^2<28\|C-I\|_\Fro^2<1$.
\end{proof}

\subsection{Deficit rules}\label{ssec:deficit}

\begin{lemma}[Deficit rules]\label{lem:deficit}
Each of the following implies \textup{(D)} on $B$.
\begin{enumerate}
\item[\textup{(i)}] A lower bound $E_-$ for $e$ on $B$ and an upper bound $D_+\le1$ for $D$ on $B$ satisfy
$E_--\eB^+-\gap^+D_+>0$. Here $D_+$ comes from Lemma~\ref{lem:Dbox} or from Lemma~\ref{lem:growth}\textup{(ii)}, and $E_-$ comes
from Lemma~\ref{lem:energy} or~\ref{lem:planarrules}, or it is the sum of a lower bound for the energy $e(\xi,f_0)$ of an anchor over
the shape box and of a lower bound for the energy gain $e(\xi,f)-e(\xi,f_0)$ over $B$ (Section~\ref{ssec:anchor};
Lemmas~\ref{lem:det-anchor} and~\ref{lem:det-centred}).
\item[\textup{(ii)}] \emph{Joint vertex rule.} Let $\mathcal T$ be a separately concave lower bound for $2\ltwo e+2\alpha+2\alpha^{-2}$
as in Lemma~\ref{lem:energy}\textup{(ii)}, let $t\in[0,1]$, and let $\mathcal C$ be a separately convex function on $B$ with
$D\le\mathcal C$ on $B$. If
\[
  \min_v\Bigl[\frac{\mathcal T(v)-2\alpha-2\alpha^{-2}}{2\ltwo}-\eB^+-\gap^+\bigl(1-t+t\,\mathcal C(v)\bigr)\Bigr]>0 ,
\]
the minimum being over the vertices $v$ of $B$, then \textup{(D)} holds on $B$.
\end{enumerate}
\end{lemma}

\begin{proof}
(i) is immediate.
(ii) $D\le1$ and $D\le\mathcal C$ give $D\le1-t+t\mathcal C$, so the bracket, a separately concave function, is a lower bound for
$\Def$, and its minimum over $B$ is attained at a vertex.
\end{proof}

For (ii) one needs a separately convex function $\mathcal C\ge D$. The certificate takes a registration $(B',f_0,n)$ and a matrix
$U\in\SO(3)$ with rational entries, and it uses the following majorant of the cost of the corresponding trial, in which
$[r_i^-,r_i^+]$ denotes the interval of $r_i$ on $B$, $i=1,2$.

\emph{The strain.} By \eqref{eq:Rchart}, the first two rows of $A=RB'^{-1}$ are affine in each of the coordinates $r_1,r_2,x,y,z$
separately, so their contribution to $\|A-U\|_\Fro^2$ is separately convex. The third row of $A$ is $w/(r_1r_2)$, where $w$ is the
third row of $B'^{-1}$, and with the third row $U_3$ of $U$ its contribution is
\[
  \Bigl|\frac w{r_1r_2}-U_3\Bigr|^2=\frac{|w|^2}{r_1^2r_2^2}-\frac{2\,w\cdot U_3}{r_1r_2}+1 ,
\]
whose second derivative with respect to $r_1$ is $(6|w|^2-4\,r_1r_2\,w\cdot U_3)/(r_1^4r_2^2)$; the same holds with $r_1$ and $r_2$
exchanged. As $|w\cdot U_3|\le|w|$, both second derivatives are nonnegative on $B$ if
\[
  3\,|w|\ \ge\ 2\,r_1^+r_2^+ ,
\]
and the rule is applied only if this holds; it is tested in the form $9|w|^2\ge4(r_1^+r_2^+)^2$, in exact rational arithmetic.

\emph{The displacement.} $|h|^2=\delta^\T R^\T R\delta$ is separately convex by the proof of Lemma~\ref{lem:sepconv}.

\emph{The unmatched neighbours.} At the points of $B$ where $q<1$, each term of $\Phi$ is, by Lemma~\ref{lem:Dbox}(ii),
$\phi(|Rv|)$ with $v=B'^{-1}b+\chi(b)\delta$ for one of the twelve vectors $b\in\Bz$ of length $\sqrt2$; it is at most
$(13-10|Rv|)_+$. Terms with $|Rv|\ge\frac{13}{10}$ on all of $B$ are omitted. For the others let $p\in\R^3$ with $|p|\le1$,
and let $\lambda_v\ge\Delta\sup_B|p_3v_3|$, with $\Delta$ as below. The verifying program uses the unit vector in the direction of
$Rv$ at the centre of $B$, which is not stored in the certificate: it checks that $Rv\ne0$ there, encloses the unit vector in a
ball, takes for $\lambda_v$ the product of upper bounds for $\Delta$ and for $|p_3v_3|$ over this ball and over $B$, and evaluates
the function $L$ below with this ball; by inclusion monotonicity the results contain the values of $L$ for the unit vector itself.
Let
$\sigma_i(t):=(r_i^-+r_i^+-t)/(r_i^-r_i^+)$ be the chord of $t\mapsto1/t$ on $[r_i^-,r_i^+]$. From
$0\le\sigma_i(t)-1/t=(t-r_i^-)(r_i^+-t)/(r_i^-r_i^+t)\le(r_i^+-r_i^-)^2/\bigl(4(r_i^-)^2r_i^+\bigr)$ and
$\sigma_1\sigma_2-\frac1{r_1r_2}=\sigma_1\bigl(\sigma_2-\frac1{r_2}\bigr)+\frac1{r_2}\bigl(\sigma_1-\frac1{r_1}\bigr)$, with
$0<\sigma_1\le1/r_1^-$ and $1/r_2\le1/r_2^-$, we get, on $B$,
\[
  0\ \le\ \sigma_1(r_1)\,\sigma_2(r_2)-\frac1{r_1r_2}\ \le\ \Delta:=\frac{(r_1^+-r_1^-)^2}{4\,(r_1^-)^2\,r_1^+\,r_2^-}+\frac{(r_2^+-r_2^-)^2}{4\,(r_2^-)^2\,r_2^+\,r_1^-} .
\]
Since $Rv=\bigl(r_1(v_1+xv_2+yv_3),\ r_2(v_2+zv_3),\ v_3/(r_1r_2)\bigr)$, the function
\[
  L:=p_1\,r_1(v_1+xv_2+yv_3)+p_2\,r_2(v_2+zv_3)+p_3\,v_3\,\sigma_1(r_1)\,\sigma_2(r_2)-\lambda_v
\]
satisfies $L\le p\cdot Rv\le|Rv|$ on $B$. It is affine in each of the eight coordinates separately, because $v$ is affine in $f$ and
does not depend on $\xi$; hence $(13-10L)_+$ is separately convex, and it is an upper bound for $\phi(|Rv|)$. At the vertices of
$B$, where $\sigma_i(r_i)=1/r_i$, the value of $L$ is $p\cdot Rv-\lambda_v$, so no chord has to be evaluated. If
$100\|A-U\|_\Fro^2+200|h|^2\le d_\sharp$ at all vertices of $B$, then $q\le d_\sharp$ on $B$ and $\Phi=0$ by
Lemma~\ref{lem:Dbox}(ii); in this case no term is added.

The function used is $\mathcal C:=100\,\|A-U\|_\Fro^2+200\,|h|^2+\sum(13-10L)_+$, the sum being over the terms that are not omitted.
It is separately convex, and $D\le\mathcal C$ on $B$: at the points where $q<1$ this follows from Lemma~\ref{lem:Dbox}(ii),
(iii) and the bounds above, and at the points where $q\ge1$ from
$\mathcal C\ge100\|A-U\|_\Fro^2+200|h|^2\ge1\ge D$, by Lemma~\ref{lem:Dbox}(iii). Rule (ii) evaluates the energy bound and the
defect bound at the same vertices; in the last stages of the cover, where the deficit is of order $10^{-10}$, this is what makes
the subdivision terminate.

\subsection{Anchor rules}\label{ssec:anchor}

Let $f_0\in\frac12\Z^3\setminus\Z^3$, $n\in\Z^3$, and for $\kappa=(\xi,f)$ let $\delta=f-f_0-n$, $h=R(\xi)\delta$ and $M=R^\T R$, so
$|h|^2=\delta^\T M\delta$. The lattice $X(\xi,f_0)$ is the \emph{anchor} of $X(\xi,f)$. Condition (A) follows from two separate
estimates, of the gain in energy and of the growth of the defect away from the anchor:
\begin{equation}\label{eq:EDg}
  e(\xi,f)-e(\xi,f_0)\ \ge\ \mathcal E(\xi,\delta),\qquad D(\xi,f)-D(\xi,f_0)\ \le\ \mathcal D(\xi,\delta),\quad\mathcal D\ge0 ,
\end{equation}
since then $\Def(\xi,f)-\Def(\xi,f_0)\ge\mathcal E-\gap^+\mathcal D$.

\begin{lemma}[Energy gain]\label{lem:gain}
\textup{(i)} $e(\xi,f_0+\delta)=e(\xi,f_0-\delta)$; in particular $\delta\mapsto e(\xi,f_0+\delta)$ has vanishing derivatives of odd
order at $0$.

\textup{(ii)} For every $\alpha>0$, with $F_{\mathrm r},F_{\mathrm d}$ as in \eqref{eq:screened},
\begin{multline}\label{eq:pairs}
  2\ltwo\bigl[e(\xi,f)-e(\xi,f_0)\bigr]\ \ge\ \sum_{\{\pm u\}}\Bigl[2F_{\mathrm r}'(q_u)\,|h|^2+4F_{\mathrm r}''\bigl(q_u+|h|^2\bigr)\,(u^\T M\delta)^2\Bigr]\\
   +\sum_{\{\pm m\}}\epsilon_m\,F_{\mathrm d}(q^*_m)\,\theta_m^2\,\Bigl(1-\frac{(1+\epsilon_m)\,\theta_m^2}{24}\Bigr)_+,
\end{multline}
where the first sum is over the pairs $\pm u$ of $\Z^3+f_0$, with $q_u=|Ru|^2$, and the second over the pairs $\pm m$ of
$\Z^3\setminus\{0\}$, with $\theta_m=2\pi m\cdot\delta$ and $\epsilon_m=-\cos(2\pi m\cdot f_0)\in\{\pm1\}$.
\end{lemma}

\begin{proof}
(i) The point reflection in the origin maps $Y(\xi,f_0+\delta)$ onto $Y(\xi,f_0-\delta)$ up to a translation by a period, because
$2f_0\in\Z^3$. (ii) In the difference the same-coset sum and the constants of \eqref{eq:screened} cancel. The vectors of the second
coset are $R(u+\delta)$, $u\in\Z^3+f_0$, and $\Z^3+f_0=-(\Z^3+f_0)$; for a pair $\pm u$ the squared lengths are
$|R(\pm u+\delta)|^2=q_u+|h|^2\pm2u^\T M\delta$. Since $F_{\mathrm r}$ is a Laplace transform of a positive measure,
$F_{\mathrm r}(Q+s)+F_{\mathrm r}(Q-s)\ge2F_{\mathrm r}(Q)+F_{\mathrm r}''(Q)s^2$ (use $\cosh\sigma\ge1+\frac12\sigma^2$ under the
integral), and $F_{\mathrm r}(q_u+|h|^2)-F_{\mathrm r}(q_u)\ge F_{\mathrm r}'(q_u)|h|^2$ by convexity. For the reciprocal pair $\pm m$ the
change is $2\epsilon_mF_{\mathrm d}(q_m^*)(1-\cos\theta_m)$, and
$\frac12\theta^2-\frac1{24}\theta^4\le1-\cos\theta\le\frac12\theta^2$, $1-\cos\theta\ge0$.
\end{proof}

To leading order in the displacement the right side of \eqref{eq:pairs} is a quadratic form in $\delta$, built from $M$, the vectors
$Mu$ and the integer vectors $m$; terms of fourth order come from the shift $|h|^2$ in the argument of $F''_{\mathrm r}$ and from the
factor $(\,\cdot\,)_+$. The certificate bounds the right side from below, for all shapes in a shape box and all displacements in a
box, by a quadratic form $\delta^\T\mathsf H(\xi)\delta$ whose matrix is enclosed uniformly over the shape box
(Lemma~\ref{lem:det-H}): the sums are truncated, with all omitted terms of unfavourable sign charged to explicit tails; the
coefficients $F'_{\mathrm r}$, $F''_{\mathrm r}$, $F_{\mathrm d}$ are replaced by affine bounds in the squared lengths, which are valid
because these functions are monotone and convex or concave; the terms of fourth order are bounded by means of an upper bound for
the displacement; and the dependence on the shape is expanded about the exact Gram matrix at the centre of the shape box, with
rigorous remainders. A second form of the energy gain, the \emph{planar gain}, uses the planar expansion \eqref{eq:planar} in the
same way, for anchors with $f_{0,3}\in\frac12+\Z$: there the difference $e(\xi,f)-e(\xi,f_0)$ is bounded below by a quadratic
expression in $\delta_3$ and $(\delta_1,\delta_2)+\delta_3\sigma$ (Lemma~\ref{lem:det-layer}). The complete statements of these
bounds are given in Appendix~\ref{app:details}.

\begin{lemma}[Growth of the defect]\label{lem:growth}
\textup{(i)} If $D(\xi,f_0)=1$, then $D(\xi,f)-D(\xi,f_0)\le0$.

\textup{(ii)} Suppose that $d_0<1$ for a registration $(B',f_0,n)$ at $\xi$. Then $D(\xi,f_0)=d_0$ and
\[
  D(\xi,f)-D(\xi,f_0)\ \le\ \min\bigl\{1-d_0,\ 200\,|h|^2+\Phi^+\bigr\},
\]
where $\Phi^+$ is an upper bound for the sum $\Phi$ of Lemma~\ref{lem:Dbox}\textup{(ii)} at $(\xi,f)$ in case $q<1$; $\Phi^+=0$ if
$q\le d_\sharp$, and in general each pair $\pm b$ of vectors of $\Bz$ of length $\sqrt2$ with $\chi(b)=1$ contributes at most
\[
  \min\Bigl\{10\,\bigl(|h|-(|Ab|-\tfrac{13}{10})\bigr)_+,\ \ 4\,\bigl(\tfrac{169}{100}-|Ab|^2+2\,|Ab\cdot h|-|h|^2\bigr)_+\Bigr\}.
\]
\end{lemma}

\begin{proof}
(i) is $D\le1$.
(ii) $D(\xi,f_0)=d_0$ by Lemma~\ref{lem:Dbox}(v), whose proof also shows that $|Ab|\ge s\sqrt2>\frac{13}{10}$ for $|b|=\sqrt2$. At
$(\xi,f)$ the same registration gives $D\le\min\{1,q+\Phi\}$, and $q-d_0=200|h|^2$; if $q\ge1$ the minimum on the right is
$1-d_0$, so let $q<1$. The vectors $Ab$ with $\chi(b)=0$ do not move. Of the two vectors $\pm Ab+h$ with $\chi(b)=1$ at most one is
shorter than $\frac{13}{10}$, because their lengths add up to at least $2|Ab|\ge\frac{13}5$; its length $\ell$ satisfies
$\ell\ge|Ab|-|h|$ and $\ell^2\ge|Ab|^2-2|Ab\cdot h|+|h|^2$, and $\ell>\frac65$ by Lemma~\ref{lem:Dbox}(ii). Its weight is
$10(\frac{13}{10}-\ell)_+$, and $10(\frac{13}{10}-\ell)\le4(\frac{169}{100}-\ell^2)$ for $\frac65\le\ell\le\frac{13}{10}$, since
$\frac{13}{10}+\ell\ge\frac52$.
\end{proof}

\begin{lemma}[Anchor rules]\label{lem:anchor}
Let $B=S\times O$, $f_0$, $n$ be as above, and let $\mathcal E$, $\mathcal D$ satisfy \eqref{eq:EDg} on $B$.
\begin{enumerate}
\item[\textup{(i)}] If $\mathcal E(\xi,\delta)-\gap^+\mathcal D(\xi,\delta)\ge c\,|h|^2$ on $B$ with $c>0$, then \textup{(A)} holds on $B$.
\item[\textup{(ii)}] If $s_-\le\Def(\xi,f_0)$ for all $\xi\in S$ and $s_-+\inf_B(\mathcal E-\gap^+\mathcal D)>0$, then \textup{(D)} holds
on $B$.
\item[\textup{(iii)}] If $D_0^-\le D(\xi,f_0)$ for all $\xi\in S$ and $\inf_B\mathcal E>\gap^+(1-D_0^-)$, then \textup{(E)} holds on $B$.
\end{enumerate}
\end{lemma}

\begin{proof}
(i) and (ii) follow from $\Def(\xi,f)-\Def(\xi,f_0)\ge\mathcal E-\gap^+\mathcal D$. (iii) By Theorem~\ref{thm:bravais},
$e(\xi,f_0)\ge\eB+\gap D_0^-$, hence $e(\xi,f)>\eB+\gap D_0^-+\gap^+(1-D_0^-)\ge\eF$.
\end{proof}

In (i) the condition is verified in most cases as the positive definiteness of a symmetric interval matrix, namely of
$\mathsf H-\gap^+\gamma\,M$ when $\mathcal E=\delta^\T\mathsf H\delta$ and $\mathcal D=\gamma|h|^2$, by an $LDL^\T$ factorization in ball
arithmetic (Lemma~\ref{lem:det-positive}); the recorded witness is then a positive pivot and not an energy margin, and the box may
contain the anchor itself. In (ii), $s_-$ may be taken to be $0$, by Theorem~\ref{thm:bravais}, or a positive lower bound
for the deficit of the anchor, obtained from a lower bound for its energy, by Lemma~\ref{lem:energy} on $S\times\{f_0+n\}$ or by
Lemma~\ref{lem:det-anchor}, and from an upper bound for its defect by Lemma~\ref{lem:Dbox}(v), with the bound $d^+$ of
Section~\ref{ssec:det-gamma} (Lemma~\ref{lem:det-centred}); a positive $s_-$ allows a form that is not positive definite. The
coefficient $\gamma$ of the defect growth is obtained from Lemma~\ref{lem:growth}: the certificate takes the largest value that
parts (i) and (ii) give over the complete list of Lemma~\ref{lem:Dbox}(iv), so that no assumption is made on which registration
is optimal (Lemma~\ref{lem:det-gamma}).

The anchor rules are the only place where Theorem~\ref{thm:bravais} enters the proof of Theorem~\ref{thm:cover}, but they are
indispensable: a neighbourhood of the family of lattices with small deficit, for instance of all lattices near the body-centred and
face-centred cubic ones, is covered by them.

\subsection{The certificate}\label{ssec:stages}

The family $\mathcal F$ of Theorem~\ref{thm:cover}, with the data that prove each of its boxes discharged, is given by a
certificate of thirteen files, which are read in order (Table~\ref{tab:stages}). At every moment there is a finite set
$\mathcal H$ of \emph{held} boxes; initially $\mathcal H=\{\Kbox\}$. One of the files marks the reflection step described below
and contains no data. Each of the other twelve replaces every box $P\in\mathcal H$ by a box $\tilde P\supset P$ together with a
subdivision of $\tilde P$ into closed boxes, called parts, each of which
\begin{itemize}
\item[(a)] is held, and then belongs to the set $\mathcal H$ of the next file; or
\item[(b)] carries condition (N), with an integer vector for which \eqref{eq:R1} or \eqref{eq:R2} is violated strictly on it, or
one of the rules of Sections~\ref{ssec:energyrules}--\ref{ssec:anchor} and Appendix~\ref{app:details}, in either case possibly
after slabs that satisfy (N) have been removed (Lemma~\ref{lem:discharged}(iii)), and with all its parameters; or
\item[(c)] is contained in $B$, in $B+(0,n)$ or in $\iota(B)+(0,n)$ with $n\in\Z^3$, for a box $B$ that the same file lists
with a rule and all its parameters.
\end{itemize}
Subdivisions are of two kinds: repeated bisection at the exact midpoint of a coordinate, both closed halves being kept;
and, in the file of stages 4 and 5, a list of boxes with rational endpoints of which it is verified in exact arithmetic
that their union contains $\tilde P$. After the last file $\mathcal H$ is empty. The following points make the
bookkeeping rigorous.

\begin{itemize}
\item \emph{Closed boxes.} Subdivisions consist of closed boxes that cover the subdivided box, and every rule is evaluated
on the closed box, the coordinate intervals being enclosed by balls. Hence, as long as no reflection has been used, the
union of the boxes of types (b), (c) and of the held boxes contains $\Kbox$ after every file, faces included. The
inequalities \eqref{eq:R1}, \eqref{eq:R2} are used for (N) only when one of them is violated strictly on the whole box for a
fixed integer vector, so no reduced point is lost at a tie.
\item \emph{Reflection.} Every box held after the third pass of stage 3 lies on one side of the hyperplane $f_1=0$. At
this point $\mathcal H$ is replaced by the set of all boxes $P\cap\iota(P')$ of positive volume with $P,P'\in\mathcal H$,
$P\subset\{f_1\ge0\}$ and $P'\subset\{f_1\le0\}$. The file of this step contains no data: the verifying program computes
the new set. The step is legitimate. Let $\mathcal Z$ be the union of the discharged boxes so far, and let
$\kappa\in\Kbox$ with $f_1\ge0$ lie neither in $\mathcal Z$ nor in $\iota(\mathcal Z)$. Then $\kappa$ has a neighbourhood
$U$ in $\Kbox\cap\{f_1\ge0\}$ with the same property. The points $\kappa'\in U$ with $f_1>0$ such that neither $\kappa'$
nor $\iota(\kappa')$ lies on a face of a held box are dense in $U$, because there are finitely many faces. Such a point
lies in the interior of a held box $P\subset\{f_1\ge0\}$, and $\iota(\kappa')$, which lies in $\Kbox\setminus\mathcal Z$
and has $f_1<0$, lies in the interior of a held box $P'\subset\{f_1\le0\}$. So $\kappa'\in P\cap\iota(P')$, and this box
has positive volume. Since the union of these finitely many boxes is closed, it contains $\kappa$. After the step,
therefore, every $\kappa\in\Kbox$ lies in a discharged box, in a held box, or in the image under $\iota$ of such a box.
This weaker statement is preserved when a held box is replaced by a subdivision as above, and when no box is held it is
the statement of Theorem~\ref{thm:cover}, because the image under $\iota$ of a discharged box is discharged. The step is
used once; the argument needs the stronger statement before it.
\item \emph{Periodicity and symmetry.} Parts of type (c) are discharged by Lemma~\ref{lem:discharged}(ii). Rules that
involve an integer vector $n$ are verified on boxes in the unwrapped offset.
\item \emph{Exact data.} The endpoints of all boxes and the data that determine the parameters of the rules are stored in the
certificate as exact rational numbers and integers: the integer vectors that violate \eqref{eq:R1} or \eqref{eq:R2}; for a
support, the value of $\alpha$, the set of terms, and either a probability distribution with dyadic weights on the vertices of a
box, of which the base points are the averages, or the instruction to take the base points at the centre or at the
vertex maxima of the box; the slabs; the registrations, the trial rotations (integer quaternions) and the numbers $t$;
the anchors. Other parameters are not stored but chosen by the verifying program from these data by fixed procedures,
which is admissible because the lemmas allow every choice; among them are the phase factors $\omega_j$ of
Lemma~\ref{lem:energy} (a lower bound for the smaller value of $w$ at the two ends of the phase interval, or $0$ if the interval
meets $\Z+\frac12$), the vectors $p$ and the numbers $\lambda_v$ of Section~\ref{ssec:deficit}, computed in ball arithmetic,
and the weights $s_T$ of Lemma~\ref{lem:energy}(iii$'$), allocated greedily up to lower bounds for $\inf_BF_j(q_j)$, the
triples being ordered in floating point. Floating-point arithmetic was used only to propose the stored data and to order these
triples.
\end{itemize}

\begin{table}[h]
\small
\begin{tabular}{l>{\raggedright\arraybackslash}p{0.66\textwidth}r}
\toprule
stage & boxes discharged in the stage, by rule & held after\\
\midrule
1 & (N): $37\,962$; energy rule (i): $28\,765$ & $966\,547$\\
2 & energy rule (ii): $243\,486$; (N): $49\,840$; planar rule (i): $43\,751$ & $675\,174$\\
3, passes 1--2 & energy rule (iii): $9\,227\,312$; energy rule (ii): $979\,538$; (N): $285\,457$; energy rule (iv): $34\,079$; planar rule (i): $27\,642$ & $597\,330$\\
3, pass 3 & energy rule (iii): $388\,387$; energy rule (ii): $220\,027$; (N): $26\,613$; energy rule (iv): $12\,247$; planar rule (i): $1\,752$; $61\,252$ of these after removing slabs & $494\,192$\\
reflection & boxes $P\cap\iota(P')$ computed by the verifier & $250\,647$\\
3, passes 4--14 & energy rule (ii): $2\,190\,441$; energy rule (iii): $1\,926\,395$; (N): $100\,898$; energy rule (iv): $53\,650$; planar rule (i): $3\,751$; anchor rule with gain (ii) $\to$ (A): $2\,302$; anchor rule (ii) with gain (ii) and $s_->0$: $289$; centred tests of Lemma~\ref{lem:det-centred} $\to$ (D): $109$; anchor rule with gain (ii) $\to$ (D): $67$; anchor rule with the planar gain $\to$ (A): $46$; centred tests of Lemma~\ref{lem:det-centred} $\to$ (E): $20$; anchor rule with the planar gain $\to$ (D): $7$; $390\,368$ of these after removing slabs & $1\,129\,620$\\
4--5 & energy rule (iii): $18\,755$; anchor rule with gain (ii) $\to$ (A): $4\,369$; anchor rule (ii) with gain (ii) and $s_->0$: $557$; deficit rule (i): $64$; anchor rule with gain (ii) $\to$ (D): $3$; contained in discharged boxes: $82\,143$ & $1\,271\,350$\\
6 & deficit rule (i): $242\,631$; energy rule (iii): $121\,658$; contained in discharged boxes: $561\,950$ & $709\,400$\\
7 & energy rule (ii): $391\,724$; deficit rule (i): $202\,119$; anchor rule with gain (ii) $\to$ (A): $82\,872$; anchor rule with the planar gain $\to$ (A): $1\,549$; anchor rule with gain (ii) $\to$ (D): $701$; anchor rule with the planar gain $\to$ (D): $230$; anchor rule with gain (ii) $\to$ (E): $1$ & $30\,258$\\
8 & deficit rule (ii): $51\,908$; energy rule (ii): $11\,435$; deficit rule (i): $11\,194$; planar rule (ii): $759$; anchor rule with the planar gain $\to$ (A): $369$ & $5\,721$\\
9 & deficit rule (ii): $42\,897$; deficit rule (i): $14\,217$; energy rule (ii): $1\,023$ & $6\,346$\\
10 & anchor rule with the planar gain $\to$ (A): $6\,324$; centred tests of Lemma~\ref{lem:det-centred} $\to$ (D): $22$ & $0$\\
\bottomrule
\end{tabular}
\caption{The certificate by stage (two files for stage 1, four for stage 3 including the reflection step, one for stages 4 and 5
together, one for each other stage) and the counts of the verification run (Section~\ref{ssec:certs}): for each stage, the boxes
discharged in it, by rule, and the number of boxes held after it. The counts by rule are those of the boxes on which the hypothesis
of a rule was checked, namely the parts of type (b) and the boxes that a file lists with a rule for parts of type (c); they add up
to the $17\,092\,214$ checks of Section~\ref{ssec:certs}(b). ``Contained in discharged boxes'' counts the parts of type (c). In
stages 4--5 and 6 every part is held or of type (c), so that the counts by rule in these rows are those of the listed boxes. Slabs
removed by Lemma~\ref{lem:discharged}(iii) are not counted separately; the counts of (N) in stage 3 include $6$ boxes (pass 3) and
$16$ boxes (passes 4--14) that satisfy (N) after slabs have been removed. Stage 1 is a dyadic partition of $\Kbox$ into $1\,033\,274$
boxes of depth up to $22$; stage 3 consists of fourteen passes, most of them adaptive bisection, with the reflection step after
its third pass. For the anchor rules, ``$\to$ (A)'', ``$\to$ (D)'' and ``$\to$ (E)'' mean parts (i), (ii) with $s_-=0$, and
(iii) of Lemma~\ref{lem:anchor}; gain (ii) is that of Lemma~\ref{lem:gain}, in the form of Lemma~\ref{lem:det-H}, and the planar
gain that of Lemma~\ref{lem:det-layer}.}\label{tab:stages}
\end{table}

\begin{proof}[Proof of Theorem~\ref{thm:cover}]
The verification (Section~\ref{ssec:certs}) establishes, in exact rational arithmetic, the statements made above about
the subdivisions and the sets $\mathcal H$, and, in ball arithmetic, the hypothesis of the rule of every part of type
(b) and of every box $B$ in (c). By Lemmas~\ref{lem:discharged}--\ref{lem:anchor} and the bounds of
Appendix~\ref{app:details}, all these boxes are discharged, that is, each satisfies (N), (E), (D) or (A), and so are the parts
of type (c), by Lemma~\ref{lem:discharged}(ii). Let $\mathcal F$ consist of all of them. By the first two points above,
every $\kappa\in\Kbox$ lies in a box of $\mathcal F$ or in the image under $\iota$ of one, since no box is held after the last
file.
\end{proof}

\begin{remark}[Zeros and smallest margins]\label{rem:margins}
The zeros of $\Def$ in $\Kbox_{\mathrm{red}}$ are the reduced representatives of the seven descriptions of each cubic
lattice as a two-site configuration (Section~\ref{sec:reduction}) and their translates by $f\mapsto f+n$. By
Theorem~\ref{thm:cover}, each of them or its
image under $\iota$ lies in a box of $\mathcal F$, and every box of $\mathcal F$ that contains a zero is discharged
by anchor rule~(i), directly or through the box that contains it (type (c)), with a pivot of at least $1.1\times10^{-5}$.
The smallest certified margin in a rule of type (E) or (D) is
$8.5\times10^{-13}$ in units of the energy per particle, to be compared with $\gap=5.6\times10^{-5}$ and with a working
precision of $128$ bits. Such a small margin reflects the crudeness of the bound that happened to be used on a particular
box: a numerical minimization of $e$ over the box in question (not certified) finds $e-\eF\approx2.7\times10^{-6}$, about
$5\%$ of $\gap$, at a vertex, and no smaller value. Where
the cover has to work hardest is near the set of lattices and near low-energy configurations with $D=1$, for instance stackings of
triangular layers, such as the hexagonal close-packed structure, whose energies are within $\gap$ of $\eF$.
\end{remark}
